\section{Introdução: o sistema sedimentar}

A sedimentologia estuda a origem, o transporte, a deposição e a transformação dos
sedimentos. Esses processos não devem ser examinados isoladamente: eles constituem
etapas conectadas do ciclo geológico e respondem às condições tectônicas, climáticas,
oceanográficas e biológicas atuantes na superfície da Terra.

De forma simplificada, o percurso que leva da rocha preexistente à formação de uma
rocha sedimentar pode ser representado por:

\[
\boxed{
\begin{gathered}
\text{rocha-fonte}
\rightarrow \text{intemperismo}
\rightarrow \text{erosão}
\rightarrow \text{transporte}
\\
\rightarrow \text{deposição}
\rightarrow \text{diagênese}
\rightarrow \text{rocha sedimentar}
\end{gathered}
}
\]

O encadeamento não é necessariamente linear. Um sedimento pode ser armazenado
temporariamente, retrabalhado várias vezes ou reincorporado ao ciclo por soerguimento e
erosão. Da mesma forma, uma rocha sedimentar pode sofrer metamorfismo, fusão ou novo
intemperismo, retornando a outras partes do ciclo das rochas.

\subsection{Controles do sistema sedimentar}

Os principais controles do sistema sedimentar incluem:

\begin{itemize}
  \item \textbf{tectônica}: produz soerguimento, subsidência, relevo e espaço de acomodação;
  \item \textbf{clima}: controla a disponibilidade de água, a temperatura e a intensidade
        do intemperismo físico e químico;
  \item \textbf{relevo}: influencia a energia potencial, as taxas de erosão e as rotas de
        transporte;
  \item \textbf{nível de base}: condiciona a erosão, o armazenamento e a deposição;
  \item \textbf{organismos}: modificam substratos, produzem sedimento e deixam estruturas
        biogênicas;
  \item \textbf{tempo}: determina quanto os processos podem alterar, selecionar e preservar
        os materiais.
\end{itemize}

\subsection{Área-fonte, rota de transporte e bacia}

Uma análise sedimentológica completa procura relacionar três domínios:

\begin{enumerate}
  \item a \textbf{área-fonte}, onde as rochas são expostas e intemperizadas;
  \item a \textbf{rota de transporte}, pela qual partículas e solutos são deslocados;
  \item a \textbf{bacia de sedimentação}, onde ocorre deposição e preservação.
\end{enumerate}

Essa abordagem permite compreender por que a composição, a granulometria, o
arredondamento, a seleção e as estruturas sedimentares registram tanto as propriedades
da área-fonte quanto a história do transporte e as condições do ambiente deposicional
\parencite{nichols2009}.
